\documentclass[12pt]{extarticle}
\newcommand{\headertext}{Week 106 / Coloring maps / Grades 2--5}
\newcommand{\packetid}{W106-grades-2-5-v1}
\usepackage[letterpaper,top=0.75in,bottom=0.8in,left=0.7in,right=0.7in,headsep=0.2in,footskip=0.4in]{geometry}
\usepackage[T1]{fontenc}
\usepackage{helvet}
\renewcommand{\familydefault}{\sfdefault}
\usepackage{tikz}
\usetikzlibrary{arrows.meta,patterns}
\usepackage{fancyhdr}
\usepackage{amssymb}
\setlength{\parindent}{0pt}
\setlength{\parskip}{0pt}
\pagestyle{fancy}
\fancyhf{}
\renewcommand{\headrulewidth}{0pt}
\lhead{\small\headertext}
\lfoot{\small Bellingham Math Circle / Week 106 / \packetid}
\rfoot{\small\thepage}
\newcommand{\fig}[1]{\input{figs/#1}}
\newcommand{\prob}[1]{\textbf{Problem #1:}~}
\newcommand{\numbox}{\tikz\draw[thick] (0,0) rectangle (1.3,1.3);}
\newcommand{\blank}{\rule{2.2cm}{0.5pt}}
\newcommand{\writelines}[1]{\par\vspace{0.3cm}\foreach \i in {1,...,#1}{\par\vspace{0.75cm}\noindent\rule{\linewidth}{0.4pt}}\par}
\newcommand{\drawbox}[2]{\tikz\draw[thick] (0,0) rectangle (#1,#2);}
\hyphenpenalty=10000
\exhyphenpenalty=10000
\raggedright

\newcommand{\umap}[1]{\begin{minipage}[t]{0.48\linewidth}\centering\fig{#1}\end{minipage}}
\newcommand{\umapn}[1]{\begin{minipage}[t]{0.48\linewidth}\centering\fig{#1}\par\vspace{0.15cm}\hfill Colors: \blank\end{minipage}}
\newcommand{\smapn}[1]{\begin{minipage}[t]{0.48\linewidth}\centering\scalebox{0.88}{\fig{#1}}\par\vspace{0.1cm}\hfill Colors: \blank\end{minipage}}
\newcommand{\dmap}[1]{\begin{minipage}[t]{0.48\linewidth}\centering\fig{#1}\par\vspace{0.2cm}\hfill Colors: \blank\end{minipage}}
\begin{document}

A map is split into countries by its lines. Each country is one piece: you can get from any spot in it to any other spot without crossing a line. The space around the shapes in a box is a country too. Countries that share a border must get different colors. Countries that meet only at a corner point may get the same color. A dot of color in a country is enough to color it.

\vspace{0.3cm}
\centerline{\fig{rule_border}\,\raisebox{0.25cm}{\LARGE$\checkmark$}\hspace{1.2cm}\fig{rule_bad}\,\raisebox{0.25cm}{\scalebox{1.4}{\LARGE$\times$}}\hspace{1.2cm}\raisebox{-0.5cm}{\fig{rule_corner}}\,\raisebox{0.25cm}{\LARGE$\checkmark$}}

\vspace{0.4cm}
\prob{1}Color each map using only two colors.

\vspace{0.4cm}
\noindent\umap{u_lines4}\hfill\umap{u_loops2}

\vspace{0.6cm}
\noindent\umap{u_venn_line}\hfill\umap{u_scribble}

\newpage
\prob{2}On your whiteboard, draw maps using only lines that run from one edge of the board to another edge, and loops that end where they started. The lines and loops may wiggle and cross each other and themselves as often as you like. Can you draw a map like this that cannot be colored with two colors? Draw the map you think is hardest to color in the box, and color it.

\vspace{0.6cm}
\centerline{\drawbox{17.5}{14}}

\newpage
\prob{3}In these maps, some lines stop partway. Color each map with as few colors as you can, and write how many colors you used. Explain why two colors are not enough for the maps where you needed more.

\vspace{0.4cm}
\noindent\smapn{u_T}\hfill\smapn{u_star6}

\vspace{0.45cm}
\noindent\smapn{u_Y}\hfill\smapn{u_dangle}

\vspace{0.45cm}
\noindent\smapn{u_ladder}\hfill\smapn{u_star5}
\writelines{3}

\newpage
\prob{4}Each map is a ring of countries around a lake. The lake and the space outside the ring are not countries. Color each ring with as few colors as you can, and write how many colors you used. How many colors would a ring of 20 countries need? A ring of 25 countries? Explain.

\vspace{0.5cm}
\noindent\umapn{u_ring4}\hfill\umapn{u_ring5}

\vspace{0.5cm}
\noindent\umapn{u_ring6}\hfill\umapn{u_ring7}
\writelines{3}

\newpage
\prob{5}In these maps, a middle country is surrounded by a ring of countries. Color each map with as few colors as you can, and write how many colors you used. How many colors would a middle country with a ring of 100 countries around it need? With a ring of 101 countries? Explain.

\vspace{0.5cm}
\noindent\umapn{u_wheel4}\hfill\umapn{u_wheel5}

\vspace{0.5cm}
\noindent\umapn{u_wheel6}\hfill\umapn{u_wheel7}
\writelines{3}

\newpage
\prob{6}Play this game with a partner and 3 colors. One of you is the Painter and the other is the Spoiler. On your turn, color any empty square with one of the 3 colors, keeping the coloring rule. A square is stuck if the squares it shares borders with already use all 3 colors. The Spoiler wins if a square gets stuck, and the Painter wins if all 9 squares get colored. Play games with each player going first. Which player can win every time? Explain how.

\vspace{0.6cm}
\noindent\fig{g_ttt}\hfill\fig{g_ttt}\hfill\fig{g_ttt}

\vspace{0.6cm}
\noindent\fig{g_ttt}\hfill\fig{g_ttt}\hfill\fig{g_ttt}
\writelines{5}

\newpage
\prob{7}Play the same game with 4 colors. A square is now stuck if the squares it shares borders with already use all 4 colors. Again play games with each player going first. Which player can win every time now? Explain how.

\vspace{0.5cm}
\noindent\fig{g_ttt}\hfill\fig{g_ttt}\hfill\fig{g_ttt}

\vspace{0.5cm}
\noindent\fig{g_ttt}\hfill\fig{g_ttt}\hfill\fig{g_ttt}
\writelines{3}

\vspace{0.7cm}
\prob{8}In Problems 1 and 2, every line runs from edge to edge or is a loop. Explain why every map made that way can be colored with two colors.
\writelines{3}

\newpage
\prob{9}Draw a map of 4 countries where each country shares a border with all three of the others. Then try to draw a map of 5 countries where each country shares a border with all four of the others.

\vspace{0.4cm}
\noindent\drawbox{8.6}{7.4}\hfill\drawbox{8.6}{7.4}

\vspace{0.9cm}
\prob{10}This map needs four colors. Can you find four of its countries that each share a border with the other three? If not, explain why there are none.

\vspace{0.3cm}
\noindent\begin{minipage}[t]{0.38\linewidth}\vspace{0pt}\scalebox{1.15}{\fig{u_wheel5_lab}}\end{minipage}\hfill
\begin{minipage}[t]{0.58\linewidth}\vspace{0pt}\writelines{6}\end{minipage}

\newpage
\prob{11}In the top box, draw a map that you think needs five colors. Your partner tries to color it with four colors by writing 1, 2, 3 or 4 in each country. Then swap jobs. Use the bottom box for another try.

\vspace{0.5cm}
\centerline{\drawbox{17.5}{10}}
\vspace{0.6cm}
\centerline{\drawbox{17.5}{10}}

\newpage
\prob{12}A donut board has its edges glued together: edge A on the left is glued to edge A on the right, and edge B at the top is glued to edge B at the bottom. Leave the board on the right and you come back on the left at the same height. Leave at the top and you come back at the bottom at the same place. Countries on opposite edges share a border along the glued edge. Each number is a different country. Color each donut map with as few colors as you can, and write how many colors you used.

\vspace{0.6cm}
\noindent\dmap{d_stripes4}\hfill\dmap{d_stripes3}

\vspace{0.7cm}
\noindent\dmap{d_grid4}\hfill\dmap{d_grid3}

\newpage
\prob{13}On a donut board, draw 5 countries where each country shares a border with all four of the others. Then try 6 countries. A country may run off one edge and come back on the glued edge, so it can show up as two or more pieces on the board: write the same number in every piece of one country.

\vspace{0.6cm}
\centerline{\fig{d_blank}}

\vspace{1cm}
\centerline{\fig{d_blank}}

\newpage
\prob{14}On this donut board, pieces with the same number join across the glued edges into one country. The second picture shows the board as the thick square, with the pieces you reach across its glued edges drawn around it. Color the board with as few colors as you can. Explain why you cannot use fewer colors.

\vspace{0.5cm}
\noindent\begin{minipage}[c]{0.5\linewidth}\centering\fig{d_seven}\end{minipage}\hfill
\begin{minipage}[c]{0.48\linewidth}\centering\fig{d_tiled7}\end{minipage}
\writelines{5}

\end{document}
